\section{线性方程为何分解而不求逆}
\label{sec:08-Numerical-Analysis/05-Numerical-Linear-Algebra/01}

一份结构分析程序，刚度矩阵 \(A\) 两千阶，载荷工况两百个：自重、风载、雪载，以及它们的各种组合。循环写得很直白，每个工况取一个右端 \(b\)，调一次 \texttt{solve(A, b)}，整夜跑完，第二天出图。后来有人提议再补三百个工况，这个循环才第一次被拿出来算账。

两百次调用里，矩阵一个数都没变。用哪一行去消哪一行、每一步的乘子是多少，两百次完全一致；变的只有跟着一起走的那个右端向量。程序把同一套消元步骤重做了两百遍，每遍算完就连同中间结果一起丢掉。

写成 \(x=A^{-1}b\) 的公式对此毫无提示，它陈述解与数据的关系，不规定这个关系该怎么兑现。关于兑现方式的工程判断，\hyperref[sec:08-Numerical-Analysis/07-Numerical-Methods-for-ML/03]{``需要的是解，而不是逆矩阵''}讲得更完整。本节问的是另一半：既然总要消元，消元本身该留下什么。

\subsection{消元的产物是一对三角矩阵}

手工解方程组时，我们先用第一行把下面各行的第一列清成零，再用第二行清第二列，直到矩阵变成上三角，然后从最后一个未知量往回代。这套动作产生了两样东西：一个上三角矩阵，就是消元结束时的形状；一堆乘子，记录了每一步用了第 \(i\) 行的多少倍去减第 \(j\) 行。把乘子按位置摆进一个单位下三角矩阵，恒等式就出来了：

\[
PA=LU,
\]

其中 \(P\) 是行置换矩阵，记录消元途中交换过哪些行；\(L\) 单位下三角，对角全为 1，其余位置装乘子；\(U\) 上三角。这个等式不是近似，在精确算术下逐元成立——它只是把消元过程重新记了一遍账。

有了它，解 \(Ax=b\) 变成先对右端做同样的置换，再走两次三角系统：

\[
Ly=Pb,\qquad Ux=y.
\]

前代与回代各约 \(n^2\) 次运算，而分解本身约 \(\tfrac23n^3\)。两者差着一整个 \(n\)。\(n=2000\) 时分解是 \(5.3\times10^9\) 次浮点运算，一次前代加回代只要 \(8\times10^6\) 次，相差六百多倍。分解一次、复用多次，就是把开头那个循环里唯一昂贵的部分从循环体里提出来。

\begin{center}
\begin{tikzpicture}[x=1cm,y=1cm,font=\small]
  \node[font=\footnotesize,black!70] at (2.3,5.6) {每个右端都从零消元};
  \node[font=\footnotesize,black!70] at (8.9,5.6) {消元一次，其余只回代};
  \foreach \y/\lb in {4.8/1,4.1/2,3.4/3,2.7/4}{
    \node[font=\footnotesize,black!70] at (-0.55,\y+0.22) {\(b_{\lb}\)};
    \draw[fill=orange!25,draw=black!45] (0,\y) rectangle (4.6,\y+0.45);
    \draw[fill=green!35,draw=black!45] (4.75,\y) rectangle (5.03,\y+0.45);
  }
  \draw[fill=orange!25,draw=black!45] (6.6,4.8) rectangle (11.2,5.25);
  \draw[fill=green!35,draw=black!45] (11.35,4.8) rectangle (11.63,5.25);
  \foreach \y in {4.1,3.4,2.7}{
    \draw[dashed,draw=black!30] (6.6,\y) rectangle (11.2,\y+0.45);
    \draw[fill=green!35,draw=black!45] (11.35,\y) rectangle (11.63,\y+0.45);
  }
  \node[font=\footnotesize] at (2.3,2.3) {\(\vdots\)};
  \node[font=\footnotesize] at (8.9,2.3) {\(\vdots\)};
  \node[font=\footnotesize] at (2.3,1.55) {合计 \(k\cdot\tfrac23n^3\)};
  \node[font=\footnotesize] at (8.9,1.55) {合计 \(\tfrac23n^3+2kn^2\)};
  \draw[fill=orange!25,draw=black!45] (0,0.45) rectangle (0.28,0.8);
  \node[right,font=\footnotesize,black!70] at (0.36,0.63) {消元 \(\tfrac23n^3\)};
  \draw[fill=green!35,draw=black!45] (3.2,0.45) rectangle (3.48,0.8);
  \node[right,font=\footnotesize,black!70] at (3.56,0.63) {前代加回代 \(2n^2\)};
  \draw[dashed,draw=black!30] (6.9,0.45) rectangle (7.18,0.8);
  \node[right,font=\footnotesize,black!70] at (7.26,0.63) {不必再做的那一遍};
\end{tikzpicture}

\emph{同一个 \(A\)，四个右端。左边把消元重做四遍，右边只做一遍，虚线框是省下来的部分。绿色那一小段若按真实比例画会细到看不见，图中已放大。}
\end{center}

回到两百个工况：左边那一列的总量是 \(1.07\times10^{12}\)，右边是 \(6.9\times10^{9}\)，一百五十多倍的差距，而且工况越多差距越大——图右侧的橙色块只有一个，往下加行只是在加绿色。同一份因子还能顺带交出别的东西：转置系统 \(A^{\trans}x=c\) 用同一对 \(L,U\) 换个次序求解；行列式是 \(U\) 对角元的连乘再配一个置换符号；LAPACK 估计条件数的例程也直接读已有的因子，不再碰原矩阵。

\subsection{一个 \(10^{-16}\) 的主元能毁掉整次消元}

消元有个前提：拿来当除数的对角元不能是零。最小的反例只要两阶，

\[
A=\begin{pmatrix}0&1\\1&1\end{pmatrix},
\]

行列式是 \(-1\)，可逆得毫无争议，但左上角摆着一个零，第一步就除不下去。交换两行，一切顺畅。

真正难对付的是接近零而不等于零的主元。把左上角换成 \(\varepsilon=10^{-16}\)，程序不会报错，它会算出乘子 \(1/\varepsilon=10^{16}\)，然后用它去缩放第一行，第二行的 \(1\) 变成 \(1-10^{16}\)。在双精度下这个数被舍入成 \(-10^{16}\)，那个原始的 \(1\) 彻底消失了。回代时 \(x_2\) 还对，\(x_1\) 却由两个 \(10^{16}\) 量级的数相减得到，有效数字所剩无几。矩阵本身条件数只有 \(2.6\)——问题一点也不病态，是算法自己制造了灾难。这正是\hyperref[sec:08-Numerical-Analysis/01-Floating-Point-and-Error/04]{``前向误差、后向误差与稳定算法''}那张四格表的左下角：良态问题上出现大误差，责任无处可推。

部分选主元的做法是，处理第 \(k\) 列时先在第 \(k\) 行及以下找到该列绝对值最大的元素，把它换到主元位置再消。这个选择直接给出一条不等式：所有乘子的绝对值都不超过 1。乘子不超过 1，消元就不会用放大的方式往矩阵里灌数，上面那种自制的抵消也就无从发生。

\subsection{部分选主元是一次有意识的权衡}

选主元有效不有效，可以用增长因子量化：

\[
\rho=\frac{\max_{i,j,k}\bigl\lvert a_{ij}^{(k)}\bigr\rvert}{\max_{i,j}\lvert a_{ij}\rvert},
\]

分子取遍消元各阶段出现过的最大元素，分母是原矩阵的最大元素。带主元 LU 的后向误差界形如 \(\norm{\Delta A}\lesssim c_n\rho\,u\norm{A}\)：机器精度打底，维数因子跟着，剩下全押在 \(\rho\) 上。

部分选主元只保证乘子不超过 1，于是每消一步元素至多翻倍，\(n\) 步下来 \(\rho\le2^{n-1}\)。这个界不是虚张声势，Wilkinson 构造过一族矩阵恰好把它顶满：对角线全 1，严格下三角全 \(-1\)，最后一列全 1。消元时最后一列会一路翻倍，\(n=60\) 就足以让双精度失去全部有效数字，而这个矩阵的条件数是 \(O(n)\)，良态得很。

所以库函数默认部分选主元，靠的不是最坏情形保证，而是三笔账加在一起划算。第一，代价几乎为零：每列扫一遍求最大值，全程 \(O(n^2)\) 次比较，与 \(\tfrac23n^3\) 相比可以忽略，而且只在列内比较，不打乱按列分块的内存访问，分块算法照样能把大部分运算压进矩阵乘法。第二，实际矩阵的 \(\rho\) 几乎总在 \(O(1)\) 到 \(O(\sqrt n)\) 之间，指数增长的例子要刻意构造才能撞上。第三，代价更高的方案换不来相称的收益：全选主元每步要在整个未处理子块里搜索最大元，比较次数升到 \(O(n^3)\)，与消元同阶，还得同时置换行与列，而它换回的增长界是 \(n^{\frac12\log n}\) 量级——理论上好得多，实践中几乎从不被需要。

这笔账的结论要连同它的前提一起记住。误差界里的 \(c_n\) 通常是 \(n\) 的低次多项式，\(n\) 上万时不再可以当成常数；\(\rho\) 是事后才知道的量，稳健的实现会在分解过程中监控它，异常时报警而不是默默返回；上面所有结论都只管后向误差，前向误差还要再乘一次 \(\kappa(A)\)，这一步谁也省不掉。

\subsection{结构决定分解还能省下多少}

到这里默认的都是一般稠密方阵，而实际问题里的矩阵很少这么一般。对称正定的矩阵不需要选主元，运算量与存储各减一半，下一节之后的 \hyperref[sec:08-Numerical-Analysis/05-Numerical-Linear-Algebra/03]{``Cholesky：利用正定结构''}专讲这件事；对称但不定的走 \(LDL^{\trans}\)，用 \(1\times1\) 与 \(2\times2\) 的对角块兼顾对称性与稳定性；方程数多于未知数的最小二乘不该走方阵这条路，那是下一节的内容。

稀疏矩阵还多一层取舍。数值上最好的主元往往落在原本为零的位置之外，选它会让 \(L\) 冒出大量新的非零元，这叫填充；只顾保持稀疏而重排行序，又可能撞上小主元。工业级求解器在数值主元与稀疏排序之间设一个阈值，允许主元不是列内最大，只要不小于最大值的某个比例即可。填充大到存不下时，就该换成只需要矩阵向量乘的迭代法，见\hyperref[sec:08-Numerical-Analysis/05-Numerical-Linear-Algebra/04]{``大型稀疏系统与 Krylov 迭代''}。

还有一种情况值得单独说：分解已经做完，才发现精度不够。此时不必重做，用原因子解一个修正方程即可。算残差 \(r_k=b-Ax_k\)，解 \(A\,\delta x_k=r_k\)，令 \(x_{k+1}=x_k+\delta x_k\)，每轮只花 \(O(n^2)\)。残差用更高精度累加时，这套迭代改进常能收回几位被舍入吞掉的数字，也支撑了低精度分解加高精度改进的混合精度做法。它的收敛条件同样写在条件数里，\(\kappa u\) 接近 1 时改进量本身也算不准。

一次分解摊销到多个右端，是解方程这件事上最基础的一次结构复用。方阵讲完，方程数多于未知数的情形要换一套几何：那里没有能精确满足的解，只有离得最近的那个点。

